\section{Discussion}
\label{sec:discussion}

\subsection{Activity, novelty and task-relevant evidence are three different things}
The empirical ordering is not the one the naive framing predicts. Raw activity is useless as a progress
signal: agents in our sample spend tens of steps busy and unproductive. Raw novelty is barely better --- an
agent that has lost the thread still reads new files and still produces new text, so ``the observation is
new'' alarms as readily on a productive window as on a stagnant one. What separates the two is whether the
new material concerns the task that was assigned. Our \evid{} channel implements that test in the cheapest
possible way: the entity is new, and its subtokens overlap the task statement more than they overlap the
rest of the corpus. It is also the largest single source of improvement over the repetition and semantic
baselines in the ablation.

That reframes the practical problem. The interesting question is not ``is the agent looping?'' but ``is the
agent learning anything about \emph{this} task?''. Looping is one symptom; irrelevant exploration is
another, and it is more common in our data.

\subsection{Why a semantic-redundancy baseline is strong but insufficient}
The livelock-detection family represents each step in an embedding space and watches for steps that stop
moving. We reproduce that idea and find it a genuinely good ranking signal (ROC-AUC \paucBFoursemantic{},
and \paucLsem{} when learned), better than repetition and much better than workspace churn. Its limitation
is structural rather than tuning-related: two consecutive steps can be maximally far apart in embedding
space while both are irrelevant to the task, which is exactly the case our labels call stagnant.
Similarity to the previous step measures \emph{information flow}, not \emph{directed} information flow. The
fix --- conditioning on the task --- is what our evidence channel adds without a second model in the loop,
which matters for deployment: a monitor that costs nothing but tokenization can run at every step of every
session, whereas a model-based judge cannot.

\subsection{The workspace channel is a trap}
Every practical monitor we know of begins with ``did the files change?''. In our labels that channel is the
weakest of the six (\paucBSevenworkspace{}), and the reason is visible in the subtypes: many stagnant
windows contain edits (edit/revert cycles, churn on the same file, edits after completion) and many
productive windows contain none (localization, hypothesis elimination, narrowing). We would keep workspace
signals in a production system --- they are cheap and catch a specific pathology --- but as one feature
among others, never the criterion.

\subsection{What a runtime should do with an alarm}
An alarm should trigger \emph{inspection}, not termination: the monitor's precision at low false-stop
budgets is high enough that a human or a supervisor model can afford to look, and the failure-predictor
literature shows that restarting is only worthwhile when the restarted trajectory gets an overlay of the
progress made so far. The alarm should carry its evidence: because our features are symbolic, an alarm can
print the entities that stopped appearing and the verification state that stopped moving, which lets an
operator decide in one glance. And budgets should be per-task rather than global, because the threshold
that is conservative on a long task is aggressive on a ten-step one.

\subsection{What we would do with more annotated data}
The limiting factor in this study is labels, not features. Two additions would change the conclusions
materially: trajectories from longer-horizon repository tasks, where stagnation can last for dozens of
steps and the annotation unit would be a genuine work episode rather than a fixed window; and an
intervention experiment --- stopping a run at an alarm and restarting it with the current diff --- which is
the only way to measure whether an alarm at a given false-stop rate produces a net improvement in task
completion rather than merely removing steps.

\section{Limitations}
\label{sec:limits}

\paragraph{Corpus.} Every quantitative claim comes from Terminal-Bench~2.0, a corpus of short,
self-contained terminal tasks (mean \pcorpusStepsMean{} steps). Our conclusions are about runs of that
length; extending them to hour-long repository work is an extrapolation, and the cross-scaffold check in
Appendix~\ref{app:transfer} is a first step rather than a substitute.

\paragraph{Annotation.} Labels were produced by language-model readers applying a written codebook, not by
human domain experts, and we state that plainly. The codebook was written before any monitor, every window
was re-read by a second reader under reduced context, and we report the resulting agreement
(\pannBinaryAgr{}\% binary, $\kappa=\pannKappa$) with where it breaks down: on window boundaries and on
partly redacted observations. Every label carries a justification citing step numbers, so the process is
auditable, but it is strong consensus annotation rather than expert adjudication. Two consequences follow.
A fixed window that cuts a productive interval in half is inherently ambiguous, which is the main source of
cross-round disagreement; and where a window contains the moment of completion, readers split and the
window is held out, so the post-completion class holds only the agreed instances of a boundary the
codebook does not define.

\paragraph{The semantic channel is a surrogate, and a real encoder did not improve it.} The intended
rolling-diversity baseline uses a pretrained sentence encoder; the model host was unreachable when the
study was run, so we substituted a hashing bag-of-words representation with the same downstream statistics.
This invalidates cross-paper numeric comparison, and our \sem{} figure should not be quoted as the state of
the art for that family. We later obtained the encoder (all-MiniLM-L6-v2) and re-ran the same four features
on the same windows and folds, changing only the representation: it scored \pprobeRocEnc{} against
\pprobeRocBoW{}. Dense embeddings make every step of a trajectory resemble every other, so rolling diversity
keeps its discriminative strength (\pprobeFeatDiversity{} two-sided) with its sign reversed, and only
nearest-neighbour similarity survives (\pprobeFeatNearest{}). Sensitivity to exact token reuse therefore
helps on this task rather than hurting, because an agent that has stopped progressing tends to reissue
identical commands.

\paragraph{The reported performance is robust, not a tuning artifact.} We tried five architecturally
distinct routes to beat the single-family baseline on the identical windows and folds: a joint model over
all six families (\pjointAuc{} against \pceilBase{}, paired difference \pjointDiff{} on
$[\pjointLo{},\pjointHi{}]$, $p=\pjointP{}$), the pretrained encoder (\pprobeRocEnc{}), multi-scale context
(0.727), a semantic relevance function (0.566) and a learned semantic model at a longer window (0.766).
None exceeds the baseline, whose task-clustered bootstrap spans $[\pceilLo{},\pceilHi{}]$ and therefore
contains every attempt. The limit appears to belong to the corpus rather than to our design choices.

\paragraph{The score is non-stationary, and that shapes the detector.} Information flow rises with run
progress --- the trend is positive in \pdrRunsPositiveTrend{} of \pdrRunsPositiveTrendMax{} runs at
\pdrMeanSlope{} per step, and the score correlates \pdrCorrStep{} with step index. Pooled thresholds and
trailing-window references therefore both fail, for opposite reasons: a pooled threshold drifts out of
calibration as the run proceeds, while a trailing reference is contaminated by the stall it is meant to
detect. The reference must be fixed early and the series smoothed before thresholding, which is what the
detector of Section~\ref{sec:res-alarm2} does. We have not established \emph{why} smoothing is sufficient,
only that it is, and the detector remains sensitive to the opening fraction: a run whose first steps are
already unproductive gives a poor reference.

\paragraph{What the release cannot support.} The natural repair for a judged label is a measured one, and
this corpus does not permit it. \pprobeRunsNoEvent{}\% of the sampled trajectories contain no verifiable
success event at all: explicit test output appears in \pprobeRunsPytest{}\% of runs and successful builds
in \pprobeRunsBuild{}\%. The one abundant marker, an exit code of zero, is anti-correlated with success ---
failed runs carry \pprobeEventsFail{} of them on average against \pprobeEventsPass{} in solved runs ---
because it records that a command did not raise, not that the task advanced. Two further weaknesses
compound this. Redaction (\pcorpusRedactFrac{}\% of observations carrying any content,
\pcorpusRedactMean{}\% averaged per trajectory) degrades precisely the evidence and verification channels
we analyse, and it is the largest single driver of annotator disagreement: fully legible windows tie
between readers at \predactTieZero{}\%, partly redacted ones at \predactTieOne{}. And relevance is computed
lexically, so a discovery described in different words than the task scores zero --- one reason the
evidence channel underperforms, and a limitation of our implementation rather than of the idea behind it.

\section{Reproducibility}
\label{sec:repro}

The pipeline is deterministic and split into stages, each writing a frozen artifact:
\code{build\_dataset.py} draws the stratified sample and writes window features;
\code{embed\_semantics.py} writes the step representations used by the semantic channel;
\code{make\_annotation\_cards.py} and its two companions write the annotation cards;
\code{build\_gold.py} merges all annotation rounds into \code{adjudicated.csv}, \code{agreement.json}
and \code{regions.json}; \code{run\_experiments.py} runs task-level cross-validation and writes
\code{window\_features.parquet} and \code{monitor\_scores.parquet} with the metric tables; and
\code{make\_figures.py} with \code{export\_results\_tex.py} regenerates every figure and every number
in this paper from those artifacts. Data are public: Terminal-Bench~2.0 trajectories
\cite{yoonholee2026tb2} and SWE-agent trajectories \cite{nebius2024sweagenttraj}. The monitor itself is
about six hundred lines of dependency-free Python (standard library plus NumPy) and runs in milliseconds
per step on one CPU core: on \prtWindows{} windows it takes \prtMsPerWindow\,ms to compute a window's
feature vector and \prtMsPerStep\,ms to score it and evaluate the alarm rule, so monitoring a
thousand-step run costs under a second. The only external service used anywhere in the pipeline is the
annotation step. A test suite (\code{run\_tests.py}) asserts the invariants the protocol depends on, in
particular that truncating a trajectory leaves the features of the surviving windows unchanged, which is
the online restriction, and a number audit (\code{audit\_numbers.py}) re-derives the headline quantities
straight from the frozen artifacts and compares them with the macros the paper prints; at the time of
writing all \pauditChecked{} checked values matched. The figures are checked the same way, by
\code{audit\_figure\_text.py}, which measures the bounding box of every drawn label and fails on any
label-on-label overlap, any label printed outside its own panel, any glyph cut at the raster border, and
any label that falls below five points once the figure is scaled to the text width.

\section{Conclusion}

We asked which observable signals let an online monitor recognise that a coding agent has stopped making
progress without knowing the correct solution. The answer our data supports is narrower and more useful
than ``detect loops''. Exact repetition is a precise but partial detector; semantic redundancy of
successive steps is a good ranking signal but blind to direction; workspace churn is close to worthless;
verification deltas are rare but nearly decisive; and the signal that adds the most is the cheapest to
compute --- whether the material the agent is currently uncovering is new \emph{and} about the task. That
finding is the one we would want a runtime builder to take away, together with the negative results: the
signals that look most natural are not the ones that work, and a monitor that stops healthy runs is worse
than no monitor at all.

\section{Author contributions and use of AI tools}

Both authors contributed to the conception of the study, the design of the codebook, the annotation
protocol, the interpretation of the results and the writing of this report. The implementation, the
annotation pipeline and the first drafts were produced with an autonomous AI coding agent; its role, the
stages it touched, the tools and versions involved and the errors it made are recorded in full in the
accompanying AI-assistance log, submitted as supporting material together with the annotation cards, the
frozen feature matrices and the monitor implementation. No result reported here was taken on trust from
that log: every number in the paper is regenerated from the frozen artifacts by the released scripts
(Section~\ref{sec:repro}), and the label set remains open to independent human re-reading.

\appendix

\section{Budget-calibrated operating points}
\label{app:calibration}
The main text reports the raw threshold sweep, in which every monitor uses the same score threshold. That
view favours monitors whose scores happen to sit on a convenient scale, so we also calibrated each
monitor's threshold on the training trajectories: for a firing budget $b$, the threshold is the $(1-b)$
quantile of that monitor's training step scores, so every monitor alarms at a comparable rate before the
persistence rule is applied. Under this calibration the ordering changes and detection collapses for most
monitors: at the $5\%$ budget the hand-designed evidence--semantic monitor catches
\pcalCThreeevidSemBFiveDet{} of annotated stagnant windows with \pcalCThreeevidSemBFiveFs{} false stops,
and the best learned semantic mixture reaches \pcalLsemBFiveDet{} with \pcalLsemBFiveFs{}. The reason is
informative rather than disappointing: the extreme scores inside a run are frequently \emph{productive}
moments --- the one decisive discovery, the test that finally passes --- so a monitor tuned to alarm only
at its own extremes is tuned to alarm in the wrong place. We therefore report both views and treat the raw
sweep as an optimistic bound.

\section{Annotation codebook}
\label{app:codebook}
The codebook given to annotators is reproduced verbatim in the released artifact
(\code{docs/annotation}\-\code{\_guide.md}). Its core is the three-channel progress test
(\S\ref{sec:definition}): a window is productive if new task-relevant information became known, or a
plausible artefact changed durably, or verification state improved; it is stagnant if none of those
happened while the agent kept acting. Six labels are available, of which two
(\textsc{blocked\_external}, \textsc{uncertain}) are held out of the binary task.

\section{Window-size sensitivity}
\label{app:windows}
The window length $w$ is the one hyperparameter the feature set has, so we re-ran the whole protocol at
$w=6$, $20$ and $30$. The window matters, and it is not monotone: at $w=6$ the monitors have too little
context and the semantic family falls to \pwSixBFoursemantic{}, with per-window detection at a $5\%$
budget about a fifth of its $w=10$ value; at $w=20$ the global ranking is similar (semantic reaches
\pwTwoZeroBFoursemantic{}) but the within-run numbers move substantially; and at $w=30$ the window is long
enough that productive work inside it masks earlier stagnation, costing the evidence channel heavily. The
$w=10$ point used throughout is thus a compromise between context and resolution rather than an optimum we
would defend on principle, and a deployment should tune it.

\section{Cross-scaffold transfer}
\label{app:transfer}
Does a monitor fitted on some agent scaffolds work on one it has never seen? We trained
the same feature sets on annotated windows from every scaffold except one and evaluated on the held-out
scaffold, for each of the nine scaffolds in the sample. The ordering from the main experiments survives,
but the numbers are lower and the spread is wide: the semantic family averages \pcssemanticMean{} (range
\pcssemanticMin{}--\pcssemanticMax{}), against \paucBFoursemantic{} trained and evaluated inside the
corpus, and every family is close to chance on at least one scaffold. Repetition (\pcsrepetitionMean{})
and workspace (\pcsworkspaceMean{}) transfer about as well as each other, and task-grounded evidence alone
transfers worst (\pcsevidenceMean{}), again improving only in combination with the semantic channel
(\pcsevidencesemanticMean{}). These signals are portable in kind but not in calibration: a deployment
should re-fit its threshold per scaffold rather than carry one across agents. A genuine cross-corpus check
--- the same extractor applied to SWE-agent trajectories over SWE-bench --- was prepared but could not be
evaluated here, because that corpus carries no stagnation annotations, and reporting alarm rates there
without labels would be a claim we cannot support. Adding that annotation set is the first item of future
work.

\section{Channel behaviour, per trajectory}
\label{app:signals}
Figure~\ref{fig:signals} in Section~\ref{sec:results} shows the three channels step by step for three
annotated trajectories, each containing both labels, so the shaded bands show where the same run stopped
progressing.
